\chapter{Discussion}\label{ch:discussion}

We formalize pattern matching in the notions of object graphs.
Although our intention is to specify patterns according to the underlying stored structures,
the nodes, edges, labels and preorder are in fact abstract, thus can represent a wider range of relations between objects.
For example, we can replace concrete attributes with getter methods. The patterns to match will not change in this case,
and the pattern graph of the object to examine can be obtained by first calling the getters.
Even if side effects exist, the semantics of evaluating the getters can be made unambiguous, possibly imitating lazy evaluation.
Edge labels can be parameter positions of a decomposition function, patterns with such positional labels can be matched by the tuple result of the decomposition.

The nodes in our pattern graphs are pure identities. Value objects have to be represented as pure nodes associated with singleton sets.
Since we treat value sets as types, matching values is classified as matching types.
Because not only an identical type matches, a subtype also matches, this gives the benefit to match a value against a range of values, for a subset can be seen as a subtype.
For example, pattern $T(e = [5,\ldots,9])$ is matched by a tree node whose element is between $5$ and $9$.
This arrangement supersedes the traditional matching of value cases.

The mechanism to specify and detect aliasing is a distinct feature of our system.
However, it also brings certain restriction. Sometimes we want to ignore the aliasing by intention, especially for immutable objects whose aliases are safe.
For example, we may write a pattern $T(c = S, d = S)$ to test if an  object has two components both of type $S$.
While an object with layout $p_3$ in \autoref{fig:junction} matches the pattern, but an object with layout $p_4$ does not.
This causes some inconvenience. In the case to ignore aliasing, we can try match conjunction $T(c = S) \gconj T(d = S)$ to split the aliased node into two component tests.
However, we cannot use match conjunction to ignore cycles, for example, the layout in \autoref{fig:obj-pattern} does not match pattern $T(l = T(r = T(l = T)))$,
because there's no way to split a node that is both a source and a target into two tests.
A better solution is to construct an intermediate object graph to split (de-alias) the nodes that are to ignore aliases and cycles beforehand,
and match the intermediate graph against the pattern.
This can be optimized into a slightly modified match algorithm, based on \autoref{algo:match-dfs-cons},
to split the nodes on the fly in construction of the match morphism.

The subtyping in our pattern matching only needs to be a preorder in the match of a single pattern.
In match conjunction, the node of an object to match must has a subtype of multiple types.
In the classical single inheritance hierarchy, unless these types form an inheritance chain, the match cannot succeed.
Modern {\OO} languages all have interfaces that support multiple inheritance, therefore, a subtype of multiple unrelated types is possible.
However, there may exist no infimum.
For a match conjunction $T(a = x: I) \gconj T(a = J)$, we may have $A$ and $B$ both extending $I$ and $J$, but the common $\inf(I, J)$ does not exist.
In this case, we must fall back the type of $x$ to $I$ even if we know that $x$ must also have type $J$ when a match succeeds.
In languages supporting intersection types~\cite{ghilezan1996strong},
$\inf(I,J) = I \cap J$ can be generated and assigned to $x$ for better typing.

On the other hand, to find the type of a variable in match disjunction is even more complicated.
Consider the following disjunction,
$$
p: T(a = I, b = J, c = x: X) \gdisj q: T(a = y: K, b = y).
$$
Not like in a dynamic language such as Python, where using undefined variables is allowed to raise run-time exceptions,
for a statically typed language, we must guarantee that a program which passes type-checking must not have such exceptions.
First, we must reject the use of any variable like $x$ which is not in the common part of the disjunction.
This is done by assigning $\NoType$ to $x$.
Second, we must collect all the types of a node that possibly occurs in any part of the disjunction with the same incoming path.
For $y$, the set of incoming paths is $\sett{[a],[b]}$, therefore, the type of $y$ must be a supertype of all the types reachable from $\sett{[a],[b]}$.
The best choice of this supertype is the supremum, $\sup(I,J,K)$, if it exists. Otherwise, we choose $\NoType$ to reject $y$,
in fact, although nondeterministic, any supertype of $I$, $J$ and $K$ will not break the type system.
Even if $y$ is properly typed, there is still a problem when an object $o$ matches $p$ rather than $q$.
Both $o.a$ and $o.b$ are legible for the initialization of $y$ with type safety.
We take nondeterministic choice here, and any implementation can decide which one to choose in a specific way.
With such complexity handling match disjunction, Scala even decides to forbid variable binding in its \emph{pattern alternatives}.

\endinput